\reportfigure{\begin{tikzpicture}
\begin{axis}[name=profile,width=.90\linewidth,height=3.8cm,xmin=-340,xmax=340,
ymin=0,ymax=1380,ytick={0,600,1200},xtick={-300,0,300},xticklabels=\empty,
ylabel={温升 (K)},ymajorgrids,legend columns=3,legend style={at={(.5,1.02)},anchor=south}]
\addplot[ink,only marks,mark=*,mark size=1.0pt] table[col sep=comma,x=figure_r_cm,y=exp_rise_K]{../../05_Calculations/op6_106_board_doorflow_90_120.csv};\addlegendentry{Peng 等}
\addplot[rust,dashed,thick] table[col sep=comma,x expr=\thisrow{model_x_cm}-330,y=fds_rise_K]{../../05_Calculations/op6_106_doorflow_90_120.csv};\addlegendentry{绝热壁}
\addplot[teal,thick] table[col sep=comma,x expr=\thisrow{model_x_cm}-330,y=fds_rise_K]{../../05_Calculations/op6_106_board_doorflow_90_120.csv};\addlegendentry{替代板材}
\end{axis}
\begin{axis}[at={(profile.south west)},anchor=north west,yshift=-.24cm,
width=.90\linewidth,height=3.0cm,xmin=-340,xmax=340,ymin=-200,ymax=200,ytick={-150,0,150},xtick={-300,0,300},xlabel={距火源轴线距离 (cm)},ylabel={误差 (K)},ymajorgrids]
\path[fill=teal!10] (axis cs:100,-200) rectangle (axis cs:330,200);
\addplot[ink!65,dashed] coordinates {(-330,0) (330,0)};
\addplot[teal,thick] table[col sep=comma,x expr=\thisrow{model_x_cm}-330,y=error_K]{../../05_Calculations/op6_106_board_doorflow_90_120.csv};
\node[font=\sffamily\scriptsize,text=ink] at (axis cs:210,125) {平均 $-112.5$ K};
\end{axis}\end{tikzpicture}}{图 2.\quad 106.2~kW、OP-6 的顶棚温升。两条 FDS 曲线采用相同诊断网格和 90--120~s 时间窗。下图为板材模型逐点预测值减去文献~\cite{peng} 的实验值的残差。}
